\chapter{Sviluppo dell'estensione}
\label{chap:sviluppo-estensione}

Questo capitolo racconta lo stage dal punto di vista del metodo di lavoro adottato, del contesto tecnico e normativo necessario a inquadrare il problema, dell'analisi dei requisiti, delle scelte di progettazione e implementazione, e infine della verifica e dei risultati raggiunti.

\section{Metodo di lavoro}
\label{sec:metodo-lavoro}

% TODO: pianificazione delle attività, modalità di interazione con il tutor aziendale,
% cadenza delle revisioni di progresso, strumenti usati per tracciare requisiti e progressi.

\section{Selezione dei criteri WCAG oggetto di verifica}
\label{sec:selezione-criteri}

% TODO: illustrazione dei tre criteri incrociati per selezionare il sottoinsieme di
% criteri WCAG 2.1 AA effettivamente implementati nell'estensione (prevalenza
% empirica, integrabilità con LLM, copertura della più ampia gamma di utenza),
% con elenco motivato dei criteri selezionati.


\section{Analisi dei requisiti}
\label{sec:analisi-requisiti}

% TODO: requisiti funzionali e non funzionali dell'estensione, casi d'uso
% principali, eventuale analisi preventiva dei rischi.


\section{Progettazione}
\label{sec:progettazione}

% TODO: architettura dell'estensione browser, scelte tecnologiche, modalità di
% integrazione con il modello linguistico, design pattern adottati.


\section{Implementazione}
\label{sec:implementazione}

% TODO: aspetti realizzativi più significativi, con frammenti di codice
% commentati dove utile.


\section{Verifica, validazione e risultati}
\label{sec:verifica-risultati}

% TODO: strategia di test ed esito; risultati qualitativi (prodotto visto dal
% lato dell'utente finale) e quantitativi (copertura requisiti, copertura test,
% metriche di codice, accuratezza della verifica automatica).
